\hspace{3mm}

\centerline{\bf \Huge 7 класс}

\hspace{3mm}

\problem{7.1} Папа три года назад был в 9 раз старше дочки. Сейчас папа в 5 раз старше дочки. Сколько лет маме, если сумма возрастов мамы, папы и дочки сейчас равна 62 года.

\problem{7.2} Два программиста написали игру на языке Python. Правила игры следующие:\par
На мониторе написано число 2023. Каждый игрок по очереди может отнимать ненулевую цифру из числа находящегося на мониторе. Побеждает тот, после чьего хода на мониторе окажется число 0. У кого из игроков есть выигрышная стратегия, вне зависимости от ходов противника? У того, кто ходит первым или вторым? \\
\problem{7.3} Лера, Аркадий и Владимир записали на доске по одному числу $-$ 0, 2 и 5 соответственно. Спустя минуту каждый из них одновременно стирает свое число и записывает вместо него сумму двух других(то есть Лера запишет 7, Аркадий 5, а Владимир 2). Спустя ещё минуту они повторяют это действие и так несколько раз. Могла ли сумма чисел Леры, Аркадия и Владимира в какой-то момент стать равной 20232022?

%регаты 7 класс 2004/2005

\problem{7.4} Докажите, что для любого натурального $n$ число $11^n + 1$ не делится на $202300$.

%https://problems.ru/view_problem_details_new.php?id=31263

\problem{7.5} Хитрый Кощей Бессмертный взял в плен в свою пещеру 21 богатыря и нарисовал на земле окружность длины 1. После чего дал богатырям испытание: «Если вы сможете встать на края окружности так, что я не найду среди вас хотя бы 100 пар богатырей, расстояние между которыми не больше $\dfrac{1}{3}$, то я вас отпущу». Докажите, что как бы богатыри не встали, им не удасться спастись от хитрого злодея. 
\par
\textit{Расстоянием} между двумя богатырями будем считать длину кратчайшей дуги от одного богатыря к другому.

%https://problems.ru/view_problem_details_new.php?id=97920